% =====================================================================
% Generative AI Assignment #1 - Technical Report
% Format: Springer LNCS
% Instructions: Paste this into Overleaf using the Springer LNCS template
%   (https://www.overleaf.com/gallery/tagged/springer)
%   Replace ROLLNO and YOUR NAME below before compiling.
% =====================================================================

\documentclass[runningheads]{llncs}

\usepackage{graphicx}
\usepackage{amsmath}
\usepackage{booktabs}
\usepackage{hyperref}
\usepackage{float}

\begin{document}

\title{Sequence Modelling for Translation, Denoising, and Generation:\\
       A Study of RNN, DAE, and VAE Architectures}

\author{Muhammad Ahmed Mufti\inst{1}}
\authorrunning{M.A. Mufti}

\institute{
    Department of Computer Science, NUCES -- FAST Islamabad\\
    Roll No: i22-1088\\
    \email{i221088@nu.edu.pk}
}

\maketitle

% =====================================================================
\begin{abstract}
This report presents implementations and evaluations of three generative
neural network architectures for distinct tasks: (1) a vanilla
Recurrent Neural Network (RNN) Encoder-Decoder for English-to-Urdu
machine translation, (2) a Convolutional Denoising Autoencoder (DAE)
for image restoration on CIFAR-10, and (3) a Variational Autoencoder
(VAE) for generative modelling on Fashion-MNIST. For translation, the
model was trained on 2,019 short sentence pairs extracted from a 9,093-
pair corpus and evaluated using BLEU score and qualitative error
analysis. The DAE was evaluated under Gaussian and salt-and-pepper
noise at multiple noise levels using PSNR and SSIM metrics, achieving
a best PSNR of 23.46\,dB. The VAE was trained across multiple latent
dimensionalities and assessed via reconstruction fidelity and the
quality of novel samples drawn from the learned prior. Key findings
confirm the well-documented limitations of vanilla RNNs for translation
and demonstrate the effectiveness of convolutional autoencoders and
VAEs for vision tasks.
\keywords{Machine Translation \and Denoising Autoencoder \and
          Variational Autoencoder \and RNN \and BLEU \and PSNR}
\end{abstract}

% =====================================================================
\section{Introduction}

Sequence-to-sequence models, autoencoders, and variational generative
models form foundational pillars of modern deep learning. This
assignment explores all three paradigms under controlled conditions:
a low-resource translation setting that stresses the limits of
unattended RNNs, a corrupted-image reconstruction task that measures
representational capacity, and a generative modelling problem that
probes the quality of a learned latent space.

The three objectives are deliberately chosen to highlight both the
strengths and failure modes of each architecture class, motivating
more sophisticated designs (attention mechanisms, diffusion models)
studied in later units.

% =====================================================================
\section{Question 1: English-to-Urdu Machine Translation}

\subsection{Dataset and Preprocessing}

The English-Urdu parallel corpus~\cite{kaggle_urdu} contains 9,102
sentence pairs. After removing 9 duplicates and filtering sentences
longer than 15 words, 2,019 usable pairs remained. These were split
80/10/10 into train (1,615), validation (201), and test (203) sets
with no overlap confirmed programmatically. Vocabulary sizes after
capping at 5,000 tokens per language were 2,438 (English) and 3,011
(Urdu). Special tokens \texttt{<pad>}, \texttt{<sos>}, \texttt{<eos>},
and \texttt{<unk>} were added to both vocabularies.

\subsection{Model Architecture}

A vanilla RNN Encoder-Decoder was implemented in PyTorch. The encoder
maps a source sequence to a fixed-size hidden vector, which initialises
the decoder. The decoder generates target tokens one step at a time
using greedy or beam-search decoding at inference.

\begin{itemize}
    \item \textbf{Encoder:} Embedding(2438, 256) $\rightarrow$ RNN(256, 512, layers=2)
    \item \textbf{Decoder:} Embedding(3011, 256) $\rightarrow$ RNN(256, 512, layers=2) $\rightarrow$ Linear(512, 3011)
    \item \textbf{Parameters:} 4,778,691
    \item \textbf{Weight init:} Orthogonal (RNN weights), zeros (biases)
\end{itemize}

\subsection{Training Setup}

Hyperparameters were selected via a grid search over four
configurations (Table~\ref{tab:q1_grid}). The best configuration
(\texttt{emb\_dim}=256, \texttt{hid\_dim}=512, \texttt{n\_layers}=2,
\texttt{lr}=0.001, \texttt{dropout}=0.3) was used for full training
over 30 epochs with Adam optimiser, teacher forcing decaying linearly
from 1.0 to 0.5, cross-entropy loss (padding tokens ignored), and
best-checkpoint saving on validation loss.

\begin{table}[H]
\caption{Hyperparameter grid search results (Q1).}
\label{tab:q1_grid}
\centering
\begin{tabular}{ccccc|c}
\toprule
\textbf{emb\_dim} & \textbf{hid\_dim} & \textbf{layers} & \textbf{lr} & \textbf{dropout} & \textbf{Val Loss}\\
\midrule
128 & 256 & 1 & 0.003 & 0.1 & 7.8065 \\
256 & 256 & 1 & 0.001 & 0.1 & 7.1212 \\
\textbf{256} & \textbf{512} & \textbf{2} & \textbf{0.001} & \textbf{0.3} & \textbf{6.8707} \\
128 & 512 & 1 & 0.003 & 0.1 & 8.4002 \\
\bottomrule
\end{tabular}
\end{table}

\subsection{Results and Evaluation}

\begin{figure}[H]
  \centering
  \includegraphics[width=0.75\textwidth]{q1_training_curves.png}
  \caption{Q1 training and validation loss over 30 epochs.}
  \label{fig:q1_curves}
\end{figure}

The training loss decreased from 6.23 to 2.92, while validation loss
increased from 6.32 to 8.03, with best checkpoint at epoch 1
(val loss 6.32). Greedy BLEU was 0.0162 and beam-search (k=3) BLEU
was 0.0180 on the test set (Table~\ref{tab:q1_bleu}).

\begin{table}[H]
\caption{Translation BLEU scores on test set (Q1).}
\label{tab:q1_bleu}
\centering
\begin{tabular}{lc}
\toprule
\textbf{Decoding Strategy} & \textbf{BLEU}\\
\midrule
Greedy & 0.0162 \\
Beam Search ($k=3$) & 0.0180 \\
\bottomrule
\end{tabular}
\end{table}

\subsection{Error Analysis}

Qualitative analysis of 30 test predictions revealed word\_repetition
in 30/30 cases — the model consistently produced high-frequency Urdu
function words (\textit{اور, نے, میں, سے}) irrespective of input.

\textbf{Root causes:}
\begin{enumerate}
    \item \textbf{Vanishing gradients:} Vanilla RNN gradients decay
    exponentially with sequence length; the decoder hidden state
    collapses to a near-constant attractor.
    \item \textbf{Fixed context bottleneck:} A single vector must
    encode the entire source sentence; without attention, the decoder
    has no mechanism to align to relevant source positions.
    \item \textbf{Severe overfitting:} With only 1,615 training pairs,
    the train/val loss gap widened monotonically (Fig.~\ref{fig:q1_curves}).
    \item \textbf{Exposure bias:} Teacher-forced training and
    autoregressive inference operate under different token distributions.
\end{enumerate}

\textbf{Future improvements:} Bahdanau attention, LSTM/GRU gating,
subword tokenisation (BPE), Transformer architecture, and a larger
parallel corpus (100k+ pairs) would each substantially improve
translation quality.

% =====================================================================
\section{Question 2: Image Denoising with Convolutional DAE}

\subsection{Dataset}

CIFAR-10~\cite{cifar10} (60,000 RGB images, 32$\times$32 pixels,
10 classes) was used. Images were normalised to $[0,1]$. Two noise
types were studied:
\begin{itemize}
    \item \textbf{Gaussian noise:} $x_{\text{noisy}} = \text{clip}(x + \mathcal{N}(0, \sigma^2), 0, 1)$
    \item \textbf{Salt-and-pepper noise:} random pixels set to 0 or 1 with probability $p$
\end{itemize}

\subsection{Model Architecture}

A symmetric convolutional DAE with configurable bottleneck channels:

\textbf{Encoder:} Conv2d(3,32) $\rightarrow$ BN $\rightarrow$ ReLU $\rightarrow$ MaxPool $\rightarrow$ Conv2d(32,64) $\rightarrow$ BN $\rightarrow$ ReLU $\rightarrow$ MaxPool\\
\textbf{Decoder:} ConvTranspose2d(64,32) $\rightarrow$ BN $\rightarrow$ ReLU $\rightarrow$ ConvTranspose2d(32,3) $\rightarrow$ Sigmoid

\subsection{Training Setup}

Both noise types were trained separately (15 epochs, Adam, lr=0.001,
MSE loss). Both training and validation MSE were tracked per epoch.
An experimental study varied noise factor across {0.1, 0.3, 0.5, 0.7}
and bottleneck channel count across {16, 32, 64, 128}.

\subsection{Results}

\begin{table}[H]
\caption{DAE reconstruction quality on test set (Q2).}
\label{tab:q2_results}
\centering
\begin{tabular}{lccc}
\toprule
\textbf{Noise Type} & \textbf{Noise Factor} & \textbf{PSNR (dB)} & \textbf{SSIM}\\
\midrule
Gaussian     & 0.3 & 21.74 & 0.944 \\
Salt-Pepper  & 0.3 & 23.46 & 0.964 \\
\bottomrule
\end{tabular}
\end{table}

\begin{figure}[H]
  \centering
  \includegraphics[width=0.85\textwidth]{q2_gaussian_results.png}
  \caption{Q2 Gaussian denoising: Original / Noisy / Reconstructed.}
  \label{fig:q2_gauss}
\end{figure}

\begin{figure}[H]
  \centering
  \includegraphics[width=0.85\textwidth]{q2_salt_pepper_results.png}
  \caption{Q2 Salt-and-pepper denoising: Original / Noisy / Reconstructed.}
  \label{fig:q2_sp}
\end{figure}

\begin{figure}[H]
  \centering
  \includegraphics[width=0.85\textwidth]{q2_training_curves.png}
  \caption{Q2 training and validation MSE loss curves for both noise types.}
  \label{fig:q2_curves}
\end{figure}

The DAE effectively removed both noise types, recovering recognisable
object structure from nearly fully corrupted inputs (middle row,
Fig.~\ref{fig:q2_gauss}). Salt-and-pepper noise achieved higher PSNR
because its pixel-level corruption is more locally predictable. Output
blurriness is characteristic of MSE-trained autoencoders, which
optimise for average pixel error.

% =====================================================================
\section{Question 3: Variational Autoencoder on Fashion-MNIST}

\subsection{Dataset}

Fashion-MNIST~\cite{fashionmnist} (70,000 grayscale images,
28$\times$28 pixels, 10 clothing categories). Pixel values normalised
to $[-1, 1]$.

\subsection{Model Architecture}

\textbf{Encoder:} FC(784$\rightarrow$512) $\rightarrow$ ReLU $\rightarrow$ FC(512$\rightarrow$256) $\rightarrow$ ReLU $\rightarrow$ FC$_\mu$(256$\rightarrow$z) + FC$_\sigma$(256$\rightarrow$z)\\
\textbf{Reparameterisation:} $z = \mu + \epsilon \cdot e^{0.5\sigma^2}$, $\epsilon \sim \mathcal{N}(0,I)$\\
\textbf{Decoder:} FC(z$\rightarrow$256) $\rightarrow$ ReLU $\rightarrow$ FC(256$\rightarrow$512) $\rightarrow$ ReLU $\rightarrow$ FC(512$\rightarrow$784) $\rightarrow$ Tanh

\subsection{Loss Function}

\begin{equation}
\mathcal{L} = \underbrace{\mathbb{E}_{q(z|x)}[\log p(x|z)]}_{\text{Reconstruction (BCE)}} - \underbrace{D_{KL}(q(z|x) \| p(z))}_{\text{Regularisation}}
\end{equation}

The KL divergence term has the closed-form solution:
$-\frac{1}{2}\sum(1 + \log\sigma^2 - \mu^2 - \sigma^2)$.

\subsection{Training Setup}

Trained 20 epochs, Adam (lr=0.001), batch size 128. An experimental
study compared latent dimensions $z \in \{2, 10, 20, 50\}$.

\subsection{Results}

\begin{figure}[H]
  \centering
  \includegraphics[width=0.85\textwidth]{q3_reconstruction.png}
  \caption{Q3 VAE: Original (top) vs Reconstructed (bottom) Fashion-MNIST images.}
  \label{fig:q3_recon}
\end{figure}

\begin{figure}[H]
  \centering
  \includegraphics[width=0.85\textwidth]{q3_generated.png}
  \caption{Q3 VAE: Novel images generated by sampling $z \sim \mathcal{N}(0,I)$.}
  \label{fig:q3_gen}
\end{figure}

\begin{figure}[H]
  \centering
  \includegraphics[width=0.85\textwidth]{q3_training_curves.png}
  \caption{Q3 VAE training curves: reconstruction loss and KL divergence.}
  \label{fig:q3_curves}
\end{figure}

Reconstructions preserved the global shape and category of each item
(shoe, trouser, shirt) while exhibiting characteristic VAE blurriness.
Novel samples drawn from the prior showed diverse, recognisable
clothing silhouettes. Higher latent dimensions ($z=50$) produced
sharper reconstructions at the cost of less smooth interpolations,
while $z=2$ collapsed category structure.

\subsection{Limitations}

VAE outputs are blurry because the pixel-wise BCE loss penalises
average error rather than perceptual quality. GANs and diffusion
models address this by replacing the reconstruction term with an
adversarial or score-matching objective.

% =====================================================================
\section{Conclusion}

This work implemented and evaluated three foundational generative
architectures. The vanilla RNN encoder-decoder confirmed theoretical
limitations: vanishing gradients, the context bottleneck, and the
inability to generalise from small corpora produce degenerate
repetitive output. These failures directly motivate attention
mechanisms and transformer architectures. The convolutional DAE
demonstrated robust denoising capability (PSNR up to 23.46\,dB,
SSIM up to 0.964), validating the encoder-bottleneck-decoder
design for image restoration. The VAE confirmed the reparameterisation
trick's effectiveness for learning a regularised latent space suitable
for generation. Future work would extend the translation system to
LSTM with Bahdanau attention, replace the DAE's MSE loss with
perceptual loss, and compare the VAE against a GAN baseline.

% =====================================================================
\begin{thebibliography}{9}

\bibitem{kaggle_urdu}
Noman, Muhammad. \textit{English to Urdu Translation Dataset}.
Kaggle, 2023. \url{https://www.kaggle.com/datasets/muhammadnoman76/translation-dataset}

\bibitem{cifar10}
Krizhevsky, A.: \textit{Learning Multiple Layers of Features from Tiny Images}.
Tech Report, University of Toronto (2009)

\bibitem{fashionmnist}
Xiao, H., Rasul, K., Vollgraf, R.: \textit{Fashion-MNIST: a Novel Image Dataset
for Benchmarking Machine Learning Algorithms}. arXiv:1708.07747 (2017)

\end{thebibliography}

\end{document}
